\scanpage{163}
On va exprimer l'axiome (iii) de façon équivalente,
en introduisant $\chi=\varphi\otimes\psi$ sur $X\boxtimes Y$.
\inserted{et tordant $\Lambda(i)$.}
On trouve alors
\deleted{sur $\operatorname{Hom}(\Lambda(-i),H^{2i}(X\boxtimes Y))$}
la forme

(iii bis) \uline{Sur} \deleted{$\unsure{H^{2i}(X)(i)|H^i(X)(i)}$}
$\operatorname{Hom}(\Lambda(-i),\unsure{\uline H^{2i}|\uline h^{2i}}(X))
\simeq\Gamma\uline h^{2i}(X)(i)$,
\uline{la forme induite par la forme} $\varphi_i$ est \uline{définie positive}.

\margin{L'équivalence de (iii) avec (iii bis) vaut moyennant (ii) et le
cas évident $(\mathrm{iii}_0)$, \unsure{utiliser} \uline{dualité} pour
exprimer les homomorphismes $X\mathbin{-}Y$ en termes d'éléments de
$\Gamma H^{2i}(X\boxtimes Y)(i)$}

Prenons $\varphi,\psi\in\Phi_X$,
\inserted{avec \deleted{$\unsure{\varphi(x,y)}$} $\widetilde\psi=\widetilde\varphi\circ A$}
Alors pour tt $u:X\longrightarrow X$, induisant
$u=u^{(i)}:\uline h^i(X)\mathbin{-}\uline h^i(X)$,
\deleted{on a \unsure{un transp.} $\unsure{u}$}
Soit $u'$ le transposé de celui-ci relatif à $\varphi,\varphi$,
alors le transposé relatif à $\varphi,\psi$ est $u''=u'A$,
et on en conclut $\operatorname{Tr}uu''\geq0$ donc $\operatorname{Tr}uu'A\geq0$.
\uline{Cela exige que} $A$ \uline{est dans le cône positif}
de l'algèbre \deleted{involutive} de Riemann
$\operatorname{End}(\uline h^i(X))$ définie par $\varphi$, --- ce qui
revient à $=$, $\varphi$ et $\psi$ sont des formes
``équivalentes'' du \unsure{pt} de vue des structures
Riemanniennes qu'elles définissent.
\mspage{162}
